\noindent
\begin{minipage}[t]{0.48\textwidth}
    \begin{center}
        BỘ GIÁO DỤC VÀ ĐÀO TẠO\\
        \textbf{ĐẠI HỌC PHENIKAA}
    \end{center}
\end{minipage}
\hfill
\begin{minipage}[t]{0.48\textwidth}
    \begin{center}
        {\fontsize{8pt}{11pt}\selectfont \textbf{CỘNG HÒA XÃ HỘI CHỦ NGHĨA VIỆT NAM}}\\
        \textbf{Độc lập - Tự do - Hạnh phúc}
    \end{center}
\end{minipage}

\vspace{0.5cm}

\begin{center}
    \textbf{\Large BIÊN BẢN ĐÁNH GIÁ ĐỒ ÁN/KHÓA LUẬN TỐT NGHIỆP}
\end{center}

\vspace{0.5cm}

% ====== I. THÔNG TIN CHUNG ======
\noindent\textbf{I. Thông tin chung (HĐ 15)}

\begin{enumerate}
    \item \textbf{Thời gian:} từ 8 giờ 30 đến 11 giờ 30 ngày 13/10/2025
    \item \textbf{Địa điểm:} Đại học Phenikaa, Tòa A6, phòng A6-203
    \item \textbf{Họ tên sinh viên:} Lê Hữu Trung \hfill \textbf{Mã SV:} 19010035
    \item \textbf{Lớp:} K13--CNTT--KHDL \hfill \textbf{Ngành:} Công nghệ thông tin \hfill \textbf{Khóa:} K13--CNTT
    \item \textbf{Tên đề tài:} \textit{Xây dựng mô hình cho bài toán người giao hàng kết hợp phương tiện bay không người lái}
    \item \textbf{Giảng viên hướng dẫn:} TS. Dương Thị Kim Huyền
\end{enumerate}

\vspace{0.4cm}
% ====== II. THÀNH PHẦN ======
\noindent\textbf{II. Thành phần}

\begin{enumerate}
    \item TS. Nguyễn Thị Thúy Liên -- Chủ tịch
    \item ThS. Nguyễn Văn Sơn -- Thư ký
    \item ThS. Phạm Trung Dũng -- Giảng viên phản biện
\end{enumerate}

\vspace{0.4cm}
% ====== III. CÂU HỎI ======
\noindent\textbf{III. Tổng hợp câu hỏi của thành viên Hội đồng}

\begin{enumerate}
    \item Chưa ước lượng được index nào tác động đến kết quả? Cái gì quan trọng nhất?
    \item Bài toán là NP-hard, hay NP-complete?
    \item Thông số biến loại nào là quan trọng? Biến nào là biến nguyên, biến nào là biến thực?
\end{enumerate}

\vspace{0.4cm}
% ====== IV. TRẢ LỜI SINH VIÊN ======
\noindent\textbf{IV. Tổng hợp nội dung trả lời của sinh viên}

\begin{enumerate}
    \item Các biến \( i, w, z \) xác định chuyến bay của drone là ảnh hưởng đến kết quả.
    \item Chưa chứng minh được trong quyển.
    \item Thời gian là biến thực; biến quyết định sử dụng xe là nhị phân; còn lại là các biến nguyên.
\end{enumerate}

\vspace{0.4cm}
% ====== V. ĐÁNH GIÁ ======
\noindent\textbf{V. Nội dung đánh giá của Hội đồng}

\begin{enumerate}
    \item \textbf{Ý nghĩa của đồ án/khóa luận:}\\
    Triển khai bài toán quy hoạch nguyên tuyến tính cho bài toán giao vận sử dụng kết hợp drone. Là bài toán có ý nghĩa khi đưa ra lời giải cơ sở cho bài toán phức tạp này.

    \item \textbf{Về nội dung, kết cấu của đồ án/khóa luận:}\\
    Nội dung mang tính hàn lâm, có trình bày rõ ràng mô hình MILP. Kết cấu đồ án đúng chuẩn một bài khóa luận tốt nghiệp.

    \item \textbf{Phương pháp nghiên cứu:}\\
    Nghiên cứu, thiết kế giải pháp 5-index, 3-index MILP.

    \item \textbf{Các kết quả nghiên cứu đạt được:}\\
    Phát triển và lập trình thành công 2 mô hình 5-index và 3-index MILP cho bài toán giao vận.

    \item \textbf{Những ưu điểm, nhược điểm và nội dung cần bổ sung, chỉnh sửa:}\\
    \textbf{Ưu điểm:}
    \begin{itemize}
        \item Không trùng lặp với các đồ án trước.
        \item Hướng nghiên cứu sâu, mức độ thực hiện có sự so sánh với các nghiên cứu trước.
    \end{itemize}
    \textbf{Nhược điểm:} Bài toán chỉ mang tính nghiên cứu, với số lượng khách hàng nhỏ (ở mức demo).
\end{enumerate}

\vspace{0.4cm}
% ====== VI. ĐIỂM KẾT LUẬN ======
\noindent\textbf{VI. Điểm kết luận:} \textit{9.4} \hfill (Bằng chữ: \textit{Chín phẩy tư})

\vspace{0.8cm}

% ====== CHỮ KÝ ======
\noindent
\begin{minipage}[t]{0.48\textwidth}
    \begin{center}
        \textbf{CHỦ TỊCH HỘI ĐỒNG}\\[0.2cm]
        (Ký, ghi rõ họ tên)\\[1cm]
        \textit{(Chữ ký)}\\[0.5cm]
        \textbf{TS. Nguyễn Thị Thúy Liên}
    \end{center}
\end{minipage}
\hfill
\begin{minipage}[t]{0.48\textwidth}
    \begin{center}
        \textbf{THƯ KÝ}\\[0.2cm]
        (Ký, ghi rõ họ tên)\\[1cm]
        \textit{(Chữ ký)}\\[0.5cm]
        \textbf{ThS. Nguyễn Văn Sơn}
    \end{center}
\end{minipage}